\documentclass[12pt, a4paper]{article}

\usepackage[french]{babel}
\usepackage[utf8]{inputenc}
\usepackage[T1]{fontenc}
\usepackage{tgadventor}
\renewcommand*\familydefault{\sfdefault}
\usepackage{amsmath, amssymb}
\usepackage{geometry}
\usepackage{enumitem}
\usepackage{multicol}
\usepackage{xcolor}
\usepackage{mdframed}
\usepackage{verbatim}

\geometry{margin=2.5cm}

\definecolor{vlmblue}{RGB}{25, 90, 160}
\definecolor{vlmgray}{RGB}{245, 245, 245}
\definecolor{vlmgreen}{RGB}{0, 130, 80}

% ── Titre de l'exercice dans un encadré ──
\newcommand{\titrexercice}[2]{%
  \begin{mdframed}[
    backgroundcolor=vlmgray,
    linecolor=vlmblue,
    linewidth=1.5pt,
    innerleftmargin=10pt,
    innerrightmargin=10pt,
    innertopmargin=6pt,
    innerbottommargin=6pt,
    skipabove=1em,
    skipbelow=0.5em
  ]
  \textbf{\textcolor{vlmblue}{Exercice #1}}%
  \ifx&#2&\else\hfill\textbf{#2}\fi
  \end{mdframed}
}

% ── Corps de l'exercice sans cadre ──
\newenvironment{exercice}[2]{%
  \titrexercice{#1}{#2}
  \vspace{0.3em}
}{%
  \vspace{1em}
}

% ── Solution ──
\ifdefined\AVECSOLUTIONS
  \newenvironment{solution}[1]{%
    \vspace{0.3em}
    \textbf{\textcolor{vlmgreen}{Solution #1}}
    \vspace{0.3em}\par
  }{%
    \vspace{1em}
  }
\else
  \newenvironment{solution}[1]{\comment}{\endcomment}
\fi

\pagestyle{empty}

\begin{document}
\begin{exercice}{EX-CH00-003}{Décomposition en facteurs premiers}

\textbf{Théorème fondamental de l'arithmétique :} \textit{« Chaque entier supérieur à 1 est premier ou un produit de nombres premiers ; la factorisation est unique, mis à part l'ordre des facteurs. »}

\medskip
Décompose chaque nombre ci-dessous en un produit de facteurs premiers.

\medskip
\begin{tabular}{p{5cm} p{5cm} p{5cm}}
\textbf{A.} & \textbf{B.} & \textbf{C.} \\[6pt]
$255  = \dotfill$ & $1340 = \dotfill$ & $720  = \dotfill$ \\[8pt]
$450  = \dotfill$ & $140  = \dotfill$ & $378  = \dotfill$ \\[8pt]
$84   = \dotfill$ & $342  = \dotfill$ & $270  = \dotfill$ \\[8pt]
$42   = \dotfill$ & $280  = \dotfill$ & $210  = \dotfill$ \\[8pt]
$600  = \dotfill$ & $144  = \dotfill$ & $117  = \dotfill$ \\[8pt]
$52   = \dotfill$ & $1450 = \dotfill$ & $480  = \dotfill$ \\[8pt]
\end{tabular}

\end{exercice}
\begin{solution}{EX-CH00-003}

\begin{tabular}{p{5cm} p{5cm} p{5cm}}
\textbf{A.} & \textbf{B.} & \textbf{C.} \\[6pt]
$255  = 3 \times 5 \times 17$          & $1340 = 2^2 \times 5 \times 67$  & $720  = 2^4 \times 3^2 \times 5$   \\[8pt]
$450  = 2 \times 3^2 \times 5^2$       & $140  = 2^2 \times 5 \times 7$   & $378  = 2 \times 3^3 \times 7$     \\[8pt]
$84   = 2^2 \times 3 \times 7$         & $342  = 2 \times 3^2 \times 19$  & $270  = 2 \times 3^3 \times 5$     \\[8pt]
$42   = 2 \times 3 \times 7$           & $280  = 2^3 \times 5 \times 7$   & $210  = 2 \times 3 \times 5 \times 7$ \\[8pt]
$600  = 2^3 \times 3 \times 5^2$       & $144  = 2^4 \times 3^2$          & $117  = 3^2 \times 13$             \\[8pt]
$52   = 2^2 \times 13$                 & $1450 = 2 \times 5^2 \times 29$  & $480  = 2^5 \times 3 \times 5$     \\[8pt]
\end{tabular}

\end{solution}
\end{document}